\Page{16}{\textenglish{04 / FUTURE}}{חדר המדידה של המחר}{42}{49}

{\fontsize{12.7}{18.1}\selectfont

על המסך, השמלה מוצגת על מישהי אחרת. אפשר לאהוב את הצבע ועדיין לא לדעת אם האורך ירגיש נכון עלייך. ביולי 2025 השיקה גוגל בארצות הברית את כלי ה־\textenglish{Try On} אחרי ניסוי מוגבל. משתמשת בוחרת פריט מתאים בממשק הקניות, מוסיפה תמונה באורך מלא ומקבלת דימוי מחולל של עצמה עם הפריט. פתאום אפשר לראות שילוב של צבע, צללית והקשר אישי בלי להתחיל אפילו בהזמנה. לפי דף העזרה של גוגל, הדימוי נועד להמחיש מראה ואינו קובע או מבטיח התאמה בפועל. \SourceRef{11}

\par}

\Photo{fitting-room.png}{77mm}{חלל מדידה עתידי מדומיין. \AIArt}

{\fontsize{12.7}{18.1}\selectfont

\StoryColumns{יש הבדל בין לראות איך לוק נראה לבין לבחור מידה. ב־\textenglish{Zalando} מתארים מערכת שמביאה בחשבון פרטי פריט מן המותג, משוב והחזרות של לקוחות וגם מדידות שמבצע צוות פנימי. החברה תיעדה בנפרד חדר מדידה וירטואלי שמאפשר להציג מידות שונות של אותו פריט על דמות תלת־ממדית לפי מידות גוף. המלצת מידה היא החלטה על פריט מסוים. החדר הווירטואלי עוזר להמחיש הבדלים. בשני המקרים, הבגדים שהקונה יכולה לבדוק תלויים בהיקף השירות ובנתוני המוצר הזמינים. \SourceRef{11}}{גם הקנייה המודרכת משתנה. אפשר להציג לקונה בגדים לפי תיאור של אירוע, סגנון או מגבלת מחיר, ולאפשר לה לצמצם את הבחירה באמצעות השיחה. גוגל הציגה חיפוש קניות ב־\textenglish{AI Mode} במסגרת ההכרזה שלה במאי 2025. סוג החיפוש הזה יכול לעזור כשמישהי יודעת מה היא רוצה להרגיש, אבל אין לה מילה מדויקת לצללית. הוא שימושי פחות אם הצעות מתאימות באמת לא מופיעות אצל הספקים שהמערכת סורקת. \SourceRef{11}}

\par}

{\fontsize{12.7}{18.1}\selectfont

החנות הדיגיטלית יכולה אפוא לתת לקוראת שלושה סוגי עזרה: להציע מה לראות, להמחיש איך אפשר ללבוש את זה ולעזור לבחור מידה. כשמפרידים ביניהם, קל יותר להחליט אם להתקדם לקנייה או לבדוק את הפריט בחנות. לפעמים כל מה שהיה חסר הוא צבע שאפשר לדמיין על הגוף. לפעמים עדיין צריך לדעת איך הבד מרגיש בישיבה.

\par}
